\subsection{The arithmetic lower bound with its exact coordinate comparison}

\begin{statement}{D.20}{Proved}{The checked Galois and full GT gamma line in every odd weight}
For every prime $\ell\ge5$ and every odd integer $m\ge3$, the
rational leading metabelian image of $G_\Q$, and therefore of
$\widehat{\rm GT}$, is exactly
\[
 \Q_\ell G_m[x,y],\qquad
 G_m=\frac{(X+Y)^m-X^m-Y^m}{XY}.
 \tag{D.20.1}
\]
The leading image is zero in weight one and every positive even
weight. In particular R9 has dimensions
$0,0,1,0,1,0,1,0,1,0$ in weights one through ten, with the
specified coordinate $p(X,Y)[x,y]$.
This verifies the lower bound used in Tasks D.1--D.3; the full
gentle and special image questions in D.1--D.2 remain separate
questions about the displayed inclusion maps.
\par\smallskip\noindent\emph{Manuscript comparison: \S51.15,
manuscript pages 725--733 in the supplied transcription,
scans 156--170. The transcription does not establish the
recto/verso labels for this range.}
\end{statement}
\begin{proof}
We use the primary formula in Ihara's ICM lecture,
\S6.2, the theorem of \S6.3, and (6.4.4)
\cite{DIharaICM1991}. In its normalization
$(f_\sigma)^{-1}=1+A_1\xi+A_2\eta$ and
$\psi_\sigma=1+A_1\xi$. For $\sigma\in\ker\chi_\ell$,
the formula is
\[
 \log\psi_\sigma^{\rm ab}
   =\sum_{\substack{r\ge3\\r\ {\rm odd}}}
       \frac{\kappa_r^{(\ell)}(\sigma)}
            {(\ell^{r-1}-1)r!}
       \big((X+Y)^r-X^r-Y^r\big),
 \quad X=\log(1+\xi),\ Y=\log(1+\eta).
 \tag{D.20.2}
\]
The cited \S6.2 states nonvanishing of these odd cyclotomic
cocycles for $\ell>2$, with value zero at complex conjugation
by their positive-real-root normalization. Formula (6.4.4)
identifies this series with the actual metabelian action.
These are the arithmetic inputs; we now supply the filtration
and coordinate arguments.

First $\kappa_r^{(\ell)}$ stays nonzero on $\ker\chi_\ell$.
Otherwise it factors through the abelian cyclotomic quotient.
The cocycle identities for a cyclotomic element $g$ and complex
conjugation $c$ would give
\[
 (1-\chi_\ell(c)^r)\kappa_r^{(\ell)}(g)
       =(1-\chi_\ell(g)^r)\kappa_r^{(\ell)}(c)=0.
\]
The left coefficient is two. This would make the whole
cocycle zero, contradicting the cited nonvanishing.

Fix $m$. Apply the actual Galois action to the rational
nilpotent completion of $F_\ell/\gamma_{m+3}F_\ell$.
We retain this extra class because a degree-$m$ multiplier
acts on the commutator in degree $m+2$.
The Hall coordinates of this finitely generated free
nilpotent pro-$\ell$ group are integral and torsion-free.
Its map to this rational completion is therefore injective,
as is the map on every lower-central graded lattice.
Consequently triviality in a specified lower degree of
the rational action detects triviality of the same
integral action; this comparison does not create a
new Galois element or weaken the required filtration.
In the normalized gauge, $x$ is fixed and the conjugator
$f$ belongs to the derived subgroup. An outer-trivial
lower term is consequently actual-trivial: its inner
generator must centralize $x$; its only possible
degree-one component is a multiple of $x$, excluded
by the normalization of $f$, and the higher-degree
centralizer is zero by the free tensor-algebra comparison.
Let $G$ be the compact image of $\ker\chi_\ell$ in its
automorphism group. It lies in the finite-dimensional
unipotent group of automorphisms acting trivially on
the abelianization. Its filtration $I^d$ is the one
defined at the start of Part D. Both the completed
group-ring action on the metabelian commutator and its
logarithm are homomorphisms on $G$: the action on its
coefficient ring is trivial when $\chi_\ell=1$.
Consequently (D.20.2) defines a nonzero additive
functional $k_m$ on $G$, proportional to $\kappa_m^{(\ell)}$.

Here is the precise Lie extraction from this compact
group. Put $\mathfrak g=\operatorname{span}_{\Q_\ell}\log G$
in the finite nilpotent automorphism Lie algebra. It is
a Lie subalgebra: for $A=\log a$, $B=\log b$, the elements
\[
 \ell^{-2n}\log[a^{\ell^n},b^{\ell^n}]
                       \longrightarrow[A,B]
\]
belong to this finite-dimensional, hence closed, span.
Choose $a_1,\ldots,a_d\in G$ whose logarithms form a
basis. The analytic map
\[
 (t_1,\ldots,t_d)\longmapsto
       \log(a_1^{t_1}\cdots a_d^{t_d})
\]
is defined for $t_i\in\Z_\ell$, takes values in
$\log G$, and has invertible differential at zero.
The inverse-function theorem therefore puts an open
neighbourhood of zero in $\mathfrak g$ inside $\log G$.
In particular every sufficiently small scalar multiple
of any vector of $\mathfrak g$ exponentiates to an
element of the actual compact image. The functional
$k_m$ extends linearly to $\mathfrak g$, either by
differentiating the displayed group-ring character or
by its additive logarithm.

Choose $\gamma\in G_\Q$ with
$u=\chi_\ell(\gamma)=1+\ell$, and write $A$ for its
conjugation operator on $\mathfrak g$. This preserves
$\mathfrak g$, because the cyclotomic kernel is normal.
On $\operatorname{gr}^d I$ it acts as multiplication by
$u^d$. Indeed its action on the free abelianization is
$u$, the graded Lie bracket multiplies degrees, and
any higher part of the automorphism changes only
higher degrees. Moreover
\[
                  k_m(AZ)=u^m k_m(Z),
\]
by the cyclotomic cocycle relation. Choose $Z$ with
$k_m(Z)\ne0$ and set
\[
                 Z_m=\prod_{d=1}^{m-1}(A-u^d)Z.
 \tag{D.20.3}
\]
Successively applying these commuting factors sends
$I^d$ into $I^{d+1}$; thus $Z_m\in I^m\mathfrak g$.
Its $k_m$ value is
$\prod_{d<m}(u^m-u^d)k_m(Z)$, which is nonzero.
There is no assertion that a vector in the ambient
target is realized merely because of its dimension:
(D.20.3) belongs to the Lie algebra of the actual
compact Galois image. The preceding open-neighbourhood
argument gives $\exp(\ell^a Z_m)\in G\cap I^m$ for
some integer $a\ge0$. It has nonzero weight-$m$
coefficient and lifts to an actual Galois element.

Finally let $h$ be this element's leading degree-$m$
conjugating Lie word and write $h=h_xx+h_yy$ in the
tensor algebra. Its commutative image is zero. The
definition of $\psi$ therefore gives, in degree $m$,
\[
 \log\psi^{\rm ab}=-(h_xx)^{\rm ab}
                      =(h_yy)^{\rm ab}.
\]
The exact Fox comparison in D.4 sends $p[x,y]$ to
$XYp$. Formula (D.20.2) thus gives a nonzero multiple
of $G_m$ in the requested coordinate. This coefficient
also survives the outer quotient. An inner derivation
in this gauge would have to centralize $x$; in degree
at least two of the metabelian free Lie algebra,
$\operatorname{ad}x$ is multiplication by $X$ on
$\Q_\ell[X,Y][x,y]$ and has zero kernel.

This proves the nonzero Galois lower bound. The actual
inclusions $G_\Q\to\widehat{\rm GT}\to
\widehat{\rm GT}_0$ carry it to the full GT image.
D.4 supplies the matching pentagon upper bound in
every odd weight and its vanishing in even weights;
weight one is zero by its explicit special equation.
Together these statements prove (D.20.1) and the table.
\end{proof}
\subsection{The full-source symmetry restriction in weight five}

\begin{statement}{D.21}{Proved}{The weight-five full special image has dimension one or three}
Let $\ell\ge5$, and let $\mathscr I_{5,\ell}$ be the rational
weight-five metabelian leading image of the full profinite group
$\operatorname{Out}^{\sharp}(\widehat F_2)$, with the source
filtration and leading map specified at the beginning of Part D.
Then
\[
             \dim_{\Q_\ell}\mathscr I_{5,\ell}\in\{1,3\}.
             \tag{D.21.1}
\]
In particular the intermediate value two offered in Task D.1 is
excluded. This statement leaves the choice between one and three
on the unresolved task ledger.
More precisely, the polynomial coordinate map
\[
 (X+Y)q(X,Y)[x,y]\longmapsto q(X,Y)
 \quad:\quad
 (X+Y)\Q_\ell[X,Y]_2[x,y]\xrightarrow{\ \sim\ }
                    \Q_\ell[X,Y]_2
 \tag{D.21.2}
\]
identifies the target with the direct sum of the gamma line
$\Q_\ell(X^2+XY+Y^2)$ and an absolutely irreducible
two-dimensional $\mathfrak S_3$-module. The image in
(D.21.1) contains that line and is stable under this actual
geometric symmetry action.
\par\smallskip\noindent\emph{Manuscript comparison: \S51.15,
manuscript pages 725--733 in the supplied transcription,
scans 156--170; the full special-image question of LM7 D.13.}
\end{statement}
\begin{proof}
We give the comparison between conjugation in the full source
and the polynomial action, including the inner corrections.
The geometric automorphisms permuting the three punctures give
an actual subgroup $\mathfrak S_3$ of
$\operatorname{Out}^{\sharp}(\widehat F_2)$. Conjugation by
this subgroup preserves every term of the source filtration:
the lower central and second-derived subgroups used to define
the target are characteristic. The leading derivation map is
equivariant because, to its least nonzero degree, conjugation
of an automorphism acts by conjugation of its derivation.
Only the degree-one part of the conjugating geometric
automorphism contributes in that degree. Thus
$\mathscr I_{5,\ell}$ is stable under the indicated
$\mathfrak S_3$ action before any representation-theoretic
argument is applied.

Put $c=[x,y]$, $X=\operatorname{ad}x$ and
$Y=\operatorname{ad}y$. Work in the free metabelian Lie
algebra, where its commutator module is
$\Q_\ell[X,Y]c$. By D.4, the exact weight-five target is
$p=(X+Y)q$, with $q$ homogeneous of degree two.
The normalized representative of its outer derivation is
\[
        D_p(x)=0,\qquad D_p(y)=[pc,y]=-Ypc.
        \tag{D.21.3}
\]
An inner derivation $\operatorname{ad}(hc)$ has coordinates
$(-Xhc,-Yhc)$. The polynomial $p$ in (D.21.3) is unique:
an inner derivation with zero $x$ coordinate has $Xh=0$,
which forces $h=0$ in the polynomial module.

Take the transposition $s:x\leftrightarrow y$ and the
three-cycle $r:x\mapsto y$, $y\mapsto-x-y$ on the
associated graded. The first sends $c$ to $-c$ and the
second fixes $c$. Their conjugated derivations have
coordinate pairs
\[
 \begin{split}
 sD_ps^{-1}&=(X p(Y,X)c,\ 0),\\
 rD_pr^{-1}&=(-(X+Y)p(Y,-X-Y)c,\ 0).
 \end{split}                                               \tag{D.21.4}
\]
Subtract respectively the inner derivations with
\[
 h_s=-p(Y,X),\qquad
 h_r=-(X+Y)q(Y,-X-Y).
\]
The second correction is a polynomial because
$p(Y,-X-Y)=-Xq(Y,-X-Y)$. Formula (D.21.3) is restored,
and the exact induced action under (D.21.2) is
\[
              s(q)=q(Y,X),\qquad r(q)=q(Y,-X-Y).
              \tag{D.21.5}
\]
The inner ambiguity has therefore been accounted for by
the displayed actual maps.

In the ordered basis $X^2,XY,Y^2$, the action matrices are
\[
 S=\begin{pmatrix}0&0&1\\0&1&0\\1&0&0\end{pmatrix},
 \qquad
 R=\begin{pmatrix}0&0&1\\0&-1&2\\1&-1&1\end{pmatrix}.
 \tag{D.21.6}
\]
They satisfy $S^2=R^3=1$ and $SRS=R^{-1}$.
Their common invariant line is generated by
$q_0=X^2+XY+Y^2$. Averaging the six matrices gives the
rank-one idempotent
\[
       P_0=\frac{1+R+R^2+S+SR+SR^2}{6}.
       \tag{D.21.7}
\]
Its kernel $U$ is a stable two-dimensional complement.
It contains no invariant vector. It also contains no
sign subrepresentation: $s(q)=-q$ forces
$q=a(X^2-Y^2)$, and $r(q)=q$ then forces $a=0$.
Every one-dimensional representation of $\mathfrak S_3$
is trivial or sign, since its abelianization has order two.
This proves that $U$ is irreducible, and the same argument
over any field extension of $\Q_\ell$ proves absolute
irreducibility. Equivalently, the sign projector is the
zero matrix on the entire three-dimensional space.

By D.20, the image $\mathscr I_{5,\ell}$ contains $G_5c$.
The exact polynomial identity
\[
           G_5=5(X+Y)(X^2+XY+Y^2)
           \tag{D.21.8}
\]
identifies this nonzero arithmetic line with $\Q_\ell q_0$.
The averaging operator $P_0$ preserves the source image,
as it is a rational linear combination of the six symmetry
operators already shown to preserve it. Hence the image
splits as that line plus its intersection with $U$.
Irreducibility makes this intersection either zero or all
of $U$. This proves (D.21.1) using the full-source
comparison; no lifting from the pro-$\ell$ target is used.

The exact program \texttt{d\_weight5\_symmetry.py} computes
both conjugation and inner-normalization maps for a general
quadratic, checks all six group matrices, both projectors,
the two-dimensional complement, and (D.21.8). Its rational
matrices and complete output are retained. The finite
certificate verifies these explicit algebraic identities;
source stability and the arithmetic inclusion are the
preceding proof and D.20.
\end{proof}

\subsection{What a finite leading congruence detects}

\begin{statement}{D.22}{Proved}{Exact integral rank recovery and a finite-congruence limitation}
Let $M$ be a finite free $\Z_\ell$-module and let $L\subset M$
be the image of a continuous homomorphism from a compact group
to the additive group of $M$. Write $L_n$ for its image under
the specified reduction $M\twoheadrightarrow M/\ell^nM$.
There are integers $e_1,\ldots,e_r\geq0$ such that
\[
 \log_\ell|L_n|=\sum_{i=1}^r\max(n-e_i,0),\qquad
 \dim_{\Q_\ell}(L\otimes\Q_\ell)
       =\lim_{n\to\infty}\frac{\log_\ell|L_n|}{n}=r.
 \tag{D.22.1}
\]
The successive difference of the displayed logarithms equals
$r$ for every $n>\max_i e_i$.

For each $a\geq1$ and each $\ell\geq5$, the actual
$\mathfrak S_3$-module $M=\Z_\ell[X,Y]_2$ of D.21 contains
two $\mathfrak S_3$-stable closed submodules with the same
image in $M/\ell^aM$, both containing $\Z_\ell q_0$, and
with rational dimensions one and three, respectively.
The corresponding two-dimensional target of D.2 contains
submodules with the same finite reduction and rational
dimensions one and two. These are explicit submodules;
we assert no realization of them as images of the profinite
sources in D.1 or D.2. They refute the inference of a
rational upper bound from one vanishing extra congruence
without further control of the integral image.
\end{statement}
\begin{proof}
Compactness makes $L$ closed. Integer multiplication preserves
this additive subgroup, and passage to limits in $M$ makes
it a $\Z_\ell$-submodule. It is finitely generated because
$\Z_\ell$ is noetherian. The Smith basis theorem gives a
basis $b_1,\ldots,b_d$ of $M$ with
\[
                 L=\bigoplus_{i=1}^r
                         \ell^{e_i}\Z_\ell b_i.
\]
The image of the $i$th summand under the specified reduction
has order $\ell^{\max(n-e_i,0)}$. Multiplying these independent
orders proves (D.22.1), including the asserted successive
difference. The leading map on the closed source filtration
in Part D is an additive continuous map; its compact image
therefore has exactly this form in an integral leading
lattice. This is the map to which the rank formula applies.

For the weight-five example, use the integral averaging
idempotent $P_0$ of D.21; its denominator six is a unit for
$\ell\geq5$. Set $U_0=\ker(P_0:M\to M)$. In the ordered
coefficient basis $X^2,XY,Y^2$, explicit bases are
\[
 q_0=(1,1,1),\qquad u_1=(1,2,0),\qquad u_2=(-1,0,1).
 \tag{D.22.2}
\]
The three-column determinant is $3$, a unit in $\Z_\ell$,
and $P_0u_1=P_0u_2=0$. Thus
$M=\Z_\ell q_0\oplus U_0$ and $U_0$ has basis $u_1,u_2$.
Both summands are stable under the exact matrices of
(D.21.6). Define
\[
 L^{(1)}=\Z_\ell q_0,\qquad
 L^{(3)}=\Z_\ell q_0\oplus\ell^a U_0.
 \tag{D.22.3}
\]
Their inclusion maps into $M$, followed by reduction modulo
$\ell^nM$, have the same image for every $n\leq a$.
Their rational ranks are one and three. Their finite-image
orders for all $n$ are, explicitly,
\[
 |L^{(1)}_n|=\ell^n,\qquad
 |L^{(3)}_n|=\ell^{\,n+2\max(n-a,0)}.
 \tag{D.22.4}
\]
This constructs the requested comparison inside the actual
symmetry module and retains both its invariant line and its
symmetry action.

For D.2 use the explicit independent weight-nine
polynomials $G_9$ and $P_9=X^2Y^2(X+Y)^3$ of D.4 and
the lattice $M_9=\Z_\ell G_9\oplus\Z_\ell P_9$ in
their rational span. The inclusions
\[
 \Z_\ell G_9\subset M_9,
 \qquad \Z_\ell G_9\oplus\ell^a\Z_\ell P_9\subset M_9
 \tag{D.22.5}
\]
give the corresponding example under
$M_9\twoheadrightarrow M_9/\ell^aM_9$.
The target is already the invariant part, so the indicated
$\mathfrak S_3$ action is trivial on both submodules.

The program \texttt{d\_leading\_lattices.py} checks the
integer matrices, the integral Smith valuations, and every
element of the finite images for $(\ell,a)=(5,2),(7,1)$
at the recorded precisions. Its complete output is retained.
The initially incorrect determinant assertion in that program
was corrected from six to the computed value three; the
failed run and its correction are retained separately.

For the suggested finite-quotient route in the task, the
actual profinite leading image maps to each finite shadow
image by reduction, as specified in D.6. An extra coordinate
vanishing in one such finite image proves divisibility of
that coordinate in the corresponding integral lattice.
Formula (D.22.3) proves that this divisibility alone does not
force its rational vanishing. A finite argument that also
controls the Smith valuations, or an integral homogeneous
identity valid at every precision, supplies additional
information; neither is silently inserted here. The integer
homogeneous pentagon identities of D.7--D.8 and the actual
arithmetic lower bound of D.20 retain their proved scopes.
\end{proof}

